\section{Supplementary Vectorization and Robustness Analyses}
\label{app:supplementary-results}

These analyses diagnose information loss after an operator program has produced
an Abstract Graph. They do not provide the primary evidence for compositional
discrimination.

\subsection{Node and Graph Vectorization}

For base node $v$ and structural bucket $j$, the node representation is
\begin{equation}
 Z_{P,b}(G)_{v,j}=\sum_{u\in V_I}
 \mathbf{1}[v\in V_{\mu(u)}]
 \mathbf{1}[\phi_b(\mu(u))=j]a(u).
\end{equation}
Graph-level features sum the node rows. With unit contributions this gives
\begin{equation}
 x_{P,b,j}(G)=\sum_{u\in V_I}|V_{\mu(u)}|
 \mathbf{1}[\phi_b(\mu(u))=j],
\end{equation}
which counts node--structure incidences rather than raw occurrences. The
supplement compares this representation with occurrence pooling obtained by
setting $a(u)=1/|V_{\mu(u)}|$.

Columns zero and one are reserved as $Z_{v,0}=1$ and
$Z_{v,1}=\deg_G(v)$. After sum pooling they equal $|V|$ and $2|E|$ under
undirected or directed total degree. All diagnostic probe results therefore
include an ablation without these columns; topology-focused synthetic pairs are
also matched on node and edge counts.

\subsection{Predictive Accessibility}

Logistic regression, ridge regression, and a linear SVM test whether retained
structural information is accessible to a simple model. Vocabularies are fitted
on the training partition, hyperparameters on validation data, and the test
partition is evaluated once. These probes diagnose vectorization and pooling;
they are not competitive downstream prediction experiments.

\begin{figure}[t]
  \centering
  \artifactfigure{fig05-predictive-accessibility.pdf}{Intrinsic discrimination
  versus linear-probe accessibility under occurrence/incidence pooling and
  reserved-feature ablations.}
  \caption{Intrinsic discrimination and predictive accessibility. This figure
  is promoted to the main paper only if it reveals a substantive and unexpected
  difference between retained information, linear accessibility, and size-based
  shortcuts.}
  \label{fig:predictive-accessibility}
\end{figure}

\subsection{Bounded Structural Identities}

With identity width $b$, structural hashes occupy $B=2^b-2$ buckets because
the first two columns are reserved. We distinguish operator collapse,
certificate equality under the declared attributed equivalence, bounded-hash
collision, and downstream distortion. Against an unbounded certificate registry
validated by exact attributed-isomorphism checks on the evaluated component
suite, the analysis reports raw collisions, frequency-weighted contamination,
linear-probe distortion, memory, and provenance ambiguity.

If certificate frequencies are $p_i$, frequency-weighted collision mass is
\begin{equation}
 C_{\mathrm{mass}}(b)=\sum_{i<j}p_i p_j
 \mathbf{1}[h_b(c_i)=h_b(c_j)].
\end{equation}
The analysis tests, rather than assumes, whether a heavy-tailed empirical
frequency distribution makes this mass smaller than the raw identity collision
rate.

\begin{figure}[t]
  \centering
  \artifactfigure{figA01-hash-width-curves.pdf}{Raw certificate collisions,
  frequency-weighted contamination, linear-probe distortion, memory, and
  provenance ambiguity across the declared identity-width grid.}
  \caption{Sensitivity to bounded structural-identity width. The unbounded
  certificate registry is validated on the evaluated component suite and used
  only as the no-bounded-hash reference.}
  \label{fig:appendix-hash-width}
\end{figure}
